\section{Realizing the local ranks in a common central sector}
\label{sec:sector}

An idempotent in a finite group algebra has a well-defined rank as a
linear operator.  We now turn each of these numerical ranks into an
isomorphism between modules over one common ring.  The enlargement of the
group is needed because the different simple or modular summands of the
original group algebra need not have the same rank fraction.

Here are the precise inputs to this step.  Let $K$ be a field of
characteristic $p$, let $M$ be a positive integer prime to $p$, and assume
that $K$ contains an element $\zeta$ of order $M$.  In the application,
$M=p^4-1$ and $|K|=p^4$. Write $C_M=\mathbb Z/M\mathbb Z$ and put
\[
 P=(\mathbb F_p)^3,\qquad A_0=(C_M)^P,\qquad
 H=A_0\rtimes P,\qquad w=|H|,
\]
where $P$ acts on the functions in $A_0$ by translation. Suppose that a
countable group $G_0$ contains specified copies $H_1,\ldots,H_m$ of $H$.
For each $i$, let $e_i\in K[H_i]$ be an idempotent, and assume that the
positive integer
\[
 d=\dim_K e_iK[H_i]
\]
is independent of $i$.  All matrices below act on column vectors in
free right modules, so their entries multiply coordinates on the left.

\begin{proposition}\label{prop:sector-ranks}
Under these assumptions, there is a group $G_1$ containing $G_0$ and a
central element $z\in G_1$ of order $M$ with the following property.  Set
\[
 R=K[G_1],\qquad
 f=\frac1M\sum_{k=0}^{M-1}\zeta^{-k}z^k,\qquad
 S=fR,\qquad P_i=fe_i.
\]
Then $f$ is a nonzero central idempotent.  There are explicitly
specified matrices
\[
 J_i\in\operatorname{Mat}_{w\times d}(S),\qquad
 Y_i\in\operatorname{Mat}_{d\times w}(S)
\]
satisfying
\begin{equation}\label{eq:sector-rectangular}
 Y_iJ_i=fI_d,\qquad J_iY_i=P_iI_w.
\end{equation}
The entries of $J_i$ and $Y_i$ lie in the group algebra of a finite
subgroup of $G_1$, multiplied by $f$.  Their construction uses only
finite group operations and linear algebra over $K$.
\end{proposition}

We prove the proposition in three stages.  First we attach the finite
groups while preserving all the given embeddings.  Next we describe
their $f$-parts as matrix algebras over group algebras of elementary
abelian $p$-groups.  Finally, nilpotence of the corresponding
augmentation ideals makes the ranks exact and produces the two
rectangular matrices.

\subsection{Attaching the finite groups}

Let $z$ initially generate an independent copy of $C_M$, and write
$Q_{\mathrm{base}}=\langle z\rangle\times G_0$.  For each $i$, define
\[
 B_i=(C_M)^{H_i},\qquad L_i=B_i\rtimes H_i,
 \qquad (h\cdot b)(x)=b(h^{-1}x).
\]
The groups $B_i$ are written additively.  Let $\delta_i\in B_i$ be the
constant function of value $1\in C_M$.  It has order $M$ and is fixed
by $H_i$, so it is central in $L_i$.  Moreover,
$\langle\delta_i\rangle\cap H_i=1$.  Thus
$C_i=\langle z\rangle\times H_i$ embeds in both $Q_{\mathrm{base}}$ and $L_i$, with
$(z^k,h)$ sent to $(k\delta_i,h)$ in the latter.  Define $G_1$ to be
the quotient of the free product
\[
 Q_{\mathrm{base}}*L_1*\cdots*L_m
\]
that identifies these two copies of $C_i$ for every $i$.

For completeness, the following action proves that all the factors
embed in this quotient.  Let $X$ be a countable disjoint union of copies
of the left regular $Q_{\mathrm{base}}$-set.  Its restriction to each $C_i$ is free and
has countably infinitely many orbits.  A countable disjoint union
$X_i$ of regular $L_i$-sets has the same property when restricted to
$C_i$.  Pair their $C_i$-orbits and choose one point in each paired
orbit; extending those point pairings equivariantly gives a
$C_i$-equivariant bijection $X\to X_i$.  Transport the regular
$L_i$-action to $X$ using this bijection.  The resulting actions of
$Q_{\mathrm{base}},L_1,\ldots,L_m$ agree on every prescribed $C_i$, hence define an
action of $G_1$ on $X$.  Each factor acts faithfully, which proves its
injectivity in $G_1$.  This argument uses the assumed embeddings
$H_i\hookrightarrow G_0$ in defining the subgroups $C_i\le Q_{\mathrm{base}}$.

The element $z$ is central in every generating factor of $G_1$ and
therefore in $G_1$; its order is still $M$ because $Q_{\mathrm{base}}$ embeds.
In $R$, direct multiplication gives
\[
 zf=\zeta f,\qquad
 f^2=\frac1M\sum_{k=0}^{M-1}\zeta^{-k}z^kf=f.
\]
The coefficient of the identity in $f$ is $M^{-1}\ne0$, since the
elements $1,z,\ldots,z^{M-1}$ are distinct.  Thus $f$ is a nonzero
central idempotent, and $S=fR$ is a ring with identity $f$.  Each $P_i$
is an idempotent in $fK[L_i]$.  Also, any identity $e_ie_j=0$ already
valid in $K[G_0]$ remains valid after multiplication by $f$ in $S$.

\subsection{The matrix blocks of the central sector}

Fix $i$ for the moment.  A character of $B_i$ means a homomorphism
$\chi:B_i\to K^*$.  If $b_x$ denotes the coordinate function with
value $1$ at $x\in H_i$ and $0$ elsewhere, then every character is
uniquely specified by exponents $a_x\in\mathbb Z/M\mathbb Z$ through
$\chi(b_x)=\zeta^{a_x}$.  Define
\[
 p_\chi=\frac1{|B_i|}\sum_{b\in B_i}\chi(b)^{-1}b
 \quad\text{in }K[B_i].
\]
Apply the Fourier calculation of Equation~\eqref{op:fourier}
to the finite index set $H_i$. The denominator is invertible because
$p\nmid M$. The elements $p_\chi$ are a basis of pairwise orthogonal
idempotents whose sum is $1$, and $bp_\chi=\chi(b)p_\chi$.

Use the action
\[
 (h\chi)(b)=\chi(h^{-1}\cdot b).
\]
Then $hp_\chi h^{-1}=p_{h\chi}$.  Since $z=\delta_i$ is fixed by
$H_i$, the set
\[
 \mathcal X_i=\{\chi:\chi(z)=\zeta\}
\]
is $H_i$-invariant, and the Fourier formulas give
\begin{equation}\label{eq:sector-fourier-f}
 f=\sum_{\chi\in\mathcal X_i}p_\chi.
\end{equation}
For each orbit in $\mathcal X_i$, choose one representative $\theta$
and put $T_\theta=\operatorname{Stab}_{H_i}(\theta)$.

\begin{lemma}\label{lem:sector-stabilizers}
For each such $\theta$, the subgroup $T_\theta$ is elementary abelian
of exponent $p$ and rank at most $3$, with the trivial subgroup allowed.
\end{lemma}

\begin{proof}
Write $\theta(b_x)=\zeta^{a_x}$.  Since
$z=\sum_{x\in H_i}b_x$, the condition on $\theta$ is
\[
 \sum_{x\in H_i}a_x=1\quad\text{in }\mathbb Z/M\mathbb Z.
\]
Suppose that a nonidentity $g\in A_0\cap T_\theta$ has order $q$.
The exponent of $A_0$ divides $M$, so $q>1$ divides $M$. Left
multiplication by $g$ partitions $H_i$ into cycles of length exactly
$q$.  Invariance of $\theta$ says that the exponents $a_x$ are
constant on these cycles.  Their sum consequently lies in
$q(\mathbb Z/M\mathbb Z)$.  It cannot equal $1$: reduction modulo
$q$ would give $0=1$.  This contradiction proves
$T_\theta\cap A_0=1$. The quotient homomorphism
$H_i\to H_i/A_0\cong P$ is therefore injective on $T_\theta$.
Every subgroup of the additive group $(\mathbb F_p)^3$ is an
$\mathbb F_p$-linear subspace, giving the assertion.
\end{proof}

Eliminating the prime-to-$p$ part of each stabilizer is what makes its
augmentation ideal nilpotent, and will permit the later change of basis.
We next give the block isomorphism, including its inverse. This is the
character-orbit construction already used in
Section~\ref{op:modular-section}, now allowing nontrivial stabilizers;
compare the Clifford-theoretic orbit picture in~\cite{Clifford1937Announcement}.
It will identify the underlying $K$-linear rank of $P_i$ in every block
with the same original integer $d$. For one orbit representative
$\theta$, abbreviate
\[
 T=T_\theta,\qquad s=[H_i:T],\qquad D=K[T].
\]
Choose representatives $r_1,\ldots,r_s$ of the left cosets $hT$.
The space $K[H_i]$ is a free right $D$-module with this ordered basis.
Let $g\in H_i$ act on it by left multiplication $L_g$, and let
$b\in B_i$ act diagonally by
\[
 D_b(h)=(h\theta)(b)h\qquad(h\in H_i).
\]
Both actions commute with the right action of $T$.  For the diagonal
action this follows from $(ht)\theta=h\theta$ for $t\in T$.
The semidirect-product relation is also respected, because
\[
 L_gD_bL_g^{-1}(h)
   =((g^{-1}h)\theta)(b)h
   =(h\theta)(g\cdot b)h
   =D_{g\cdot b}(h).
\]
Thus we have an algebra homomorphism
$\rho_\theta:K[L_i]\to\operatorname{Mat}_s(D)$.
The element $z$ acts as $\zeta$, so $\rho_\theta(f)=I_s$.

For $t_0\in T$ and $1\le a,b\le s$, consider
\begin{equation}\label{eq:sector-block-inverse}
 p_{r_a\theta}\,r_at_0r_b^{-1}.
\end{equation}
It belongs to $fK[L_i]$ by \eqref{eq:sector-fourier-f}.  In the
$\theta$-representation it sends $r_bt'$ to $r_at_0t'$ for every
$t'\in T$ and vanishes on every other basis coset.  To check the last
assertion, its left group factor sends the coset $r_cT$ to
$r_at_0r_b^{-1}r_cT$; the Fourier projection keeps this coset precisely
when it equals $r_aT$, which is equivalent to $c=b$.  In a
representation associated to another orbit, the Fourier projection
vanishes identically.  Consequently \eqref{eq:sector-block-inverse}
maps to the matrix with entry $t_0$ in position $(a,b)$ of the
$\theta$-block and zero elsewhere.

Taking the product of these representations gives a surjective map
\begin{equation}\label{eq:sector-blocks}
 \rho:fK[L_i]\longrightarrow
 \prod_{\theta}\operatorname{Mat}_{[H_i:T_\theta]}(K[T_\theta]).
\end{equation}
It is injective as well.  The elements $p_\chi h$, for
$\chi\in\mathcal X_i$ and $h\in H_i$, form a basis of its source.
An orbit of size $s$ contributes $sw=s^2|T|$ basis elements, equal to
the dimension of the corresponding target block.  Hence the source
and target have equal finite $K$-dimension.  Formula
\eqref{eq:sector-block-inverse}, extended linearly in $t_0$ and summed
over the positions and blocks, is an explicit inverse of $\rho$.

\subsection{Turning dimensions into rectangular identities}

Continue in one block and let $E=\rho_\theta(P_i)\in
\operatorname{Mat}_s(D)$.  As a $K$-linear operator on $K[H_i]$,
$E$ is exactly left multiplication by $e_i$, since $f$ acts as the
identity.  Therefore
\begin{equation}\label{eq:sector-regular-rank}
 \dim_K\operatorname{im}E=d.
\end{equation}
This is the link between the original rank calculation and the block
decomposition.  We will now obtain a basis in which $E$ is a coordinate
projection over $D$, not just over $K$. The argument is a constructive
nilpotent-ideal instance of the freeness of projective modules over local
rings, studied by Kaplansky~\cite[Theorem~2]{Kaplansky1958}; no general projective-module
theorem is needed here.

Let $\varepsilon:D\to K$ be augmentation, and set
$\mathfrak a=\ker\varepsilon$.  By Lemma
\ref{lem:sector-stabilizers}, we may choose generators
$t_1,\ldots,t_q$ of $T\cong(C_p)^q$ with $0\le q\le3$.
The elements $X_j=t_j-1$ commute, satisfy $X_j^p=0$, and generate
$\mathfrak a$.  Expanding $t_j=1+X_j$ and counting dimensions gives
\[
 D\cong K[X_1,\ldots,X_q]/(X_1^p,\ldots,X_q^p).
\]
It follows that
\[
 \mathfrak a^{D_{\max}+1}=0,\qquad D_{\max}=3(p-1).
\]
Indeed, a monomial of total degree greater than $q(p-1)$ has an
exponent at least $p$.  This also covers $q=0$, when
$D=K$ and $\mathfrak a=0$.  Matrices with entries in $\mathfrak a$
have the same nilpotence bound under multiplication, because each
entry of a product of $D_{\max}+1$ such matrices lies in
$\mathfrak a^{D_{\max}+1}$.

Put $\overline E=\varepsilon(E)$, applying augmentation entrywise,
and let $d_0=\operatorname{rank}_K\overline E$.  As
$\overline E$ is idempotent, its image and the image of
$I_s-\overline E$ form complementary subspaces of $K^s$.
Scan the columns of $\overline E$ in order and keep each column that
increases their span; then do the same for $I_s-\overline E$.
Build $Q$ by taking the corresponding $d_0$ columns of $E$, followed
by the corresponding $s-d_0$ columns of $I_s-E$.  Its augmentation
$Q_0$ is invertible over $K$.  Regard $Q_0$ as a scalar matrix over
$D$, and set $N=Q_0^{-1}(Q-Q_0)$.  Its entries belong to
$\mathfrak a$, so $N^{D_{\max}+1}=0$.  Since $Q=Q_0(I_s+N)$,
\begin{equation}\label{eq:sector-Q-inverse}
 Q^{-1}=\left(\sum_{k=0}^{D_{\max}}(-N)^k\right)Q_0^{-1}.
\end{equation}
Each of the first $d_0$ columns of $Q$ is fixed by $E$, and every
remaining column is killed by $E$.  Thus, for
$\Pi=\operatorname{diag}(I_{d_0},0)$,
\[
 EQ=Q\Pi,\qquad Q^{-1}EQ=\Pi.
\]
Conjugation by $Q$ preserves $K$-dimension of images.  Comparing with
\eqref{eq:sector-regular-rank} gives the integer equalities
\begin{equation}\label{eq:sector-dimensions}
 d=d_0|T|,\qquad w=s|T|,\qquad wd_0=sd.
\end{equation}
In particular, these are equalities of ordinary dimensions, not
equalities after reduction to $K$.

We can now match coordinates exactly.  Order pairs lexicographically
and choose the increasing bijection between
\[
 \{(\mu,l):1\le\mu\le d,\ 1\le l\le s\}
 \quad\text{and}\quad
 \{(\lambda,k):1\le\lambda\le w,\ 1\le k\le d_0\};
\]
the sets have equal size by \eqref{eq:sector-dimensions}.
Let $V$ be the $w$-by-$d$ block matrix with $s$-by-$s$ entries
that sends coordinate $(\mu,l)$ to its paired coordinate
$(\lambda,k)$.  Thus $V_{\lambda\mu}$ is the sum of standard
matrix units $E_{kl}^{\mathrm{std}}$ over all such pairings.
Let $W$ be the $d$-by-$w$ block matrix reversing each pairing, with
$W_{\mu\lambda}$ the corresponding sum of
$E_{lk}^{\mathrm{std}}$.  It sends all unpaired coordinates with
$k>d_0$ to zero.  Directly on these coordinates,
\[
 WV=I_d\otimes I_s,\qquad VW=I_w\otimes\Pi.
\]
Define block matrices
\[
 J^{(\theta)}=(I_w\otimes Q)V,\qquad
 Y^{(\theta)}=W(I_w\otimes Q^{-1}).
\]
They have outer shapes $w\times d$ and $d\times w$, respectively,
with entries in $\operatorname{Mat}_s(D)$.  Their products are
\[
 Y^{(\theta)}J^{(\theta)}=I_d\otimes I_s,\qquad
 J^{(\theta)}Y^{(\theta)}=I_w\otimes E.
\]
Finally define the $(\lambda,\mu)$ entry of $J_i$ by applying
$\rho^{-1}$ to the tuple
$(J^{(\theta)}_{\lambda\mu})_\theta$, and similarly define $Y_i$
entrywise from the $Y^{(\theta)}$.  The isomorphism
\eqref{eq:sector-blocks} carries $f$ to the tuple of identity
matrices and $P_i$ to the tuple of matrices $E$.  The last two
identities therefore pull back to
\eqref{eq:sector-rectangular}.  This proves Proposition
\ref{prop:sector-ranks}.

All choices in the matrix construction are over finite sets: character orbits, coset
representatives, augmentation matrices, and coordinate pairings.
Taking the first element or the first independent columns in fixed
orders makes the prescription definite.  The inverse
\eqref{eq:sector-Q-inverse} is a finite sum, and the inverse
\eqref{eq:sector-block-inverse} is a finite group-algebra formula.
Thus the matrices furnished by the proposition are finite sums in
the explicitly attached finite groups.

We now apply Proposition~\ref{prop:sector-ranks} to the embedded groups
of Proposition~\ref{prop:first-gluing-injective} and the identical
idempotents of Proposition~\ref{op:modular}. The group $G_0$
is countable because it is finitely generated. All the ranks equal the
positive integer $d$ from Equation~\eqref{op:regular-rank}.
Hence the group $G_1$, the nonzero central idempotent $f$, and the
matrices $J_i,Y_i$ are now fixed by the constructions above.
Corollary~\ref{op:edge-zero} also gives $P_iP_j=0$ on every
edge with $i<j$. These are precisely the inputs to the next section.
